\section{Discussion}
\label{sec:discussion}

This study is designed to distinguish three questions that a single
Dijkstra-versus-policy comparison cannot answer: whether heuristic pruning
changes the classical compute baseline, whether incremental reuse changes it
again, and whether RL behavior is stable across training seeds and distribution
shifts. The generalization, scaling, realism, and ablation tables should
therefore be interpreted together rather than reducing performance to one
headline average.

Classical methods retain an important structural advantage: on the currently
observed graph they provide an optimal route, and failures arise from future
unobserved changes or route semantics rather than approximation by the search
itself. A* measures how far simple heuristic pruning goes; D* Lite measures the
additional value of reusing prior search effort. The measured data do not show
a deployment-time crossover in favor of RL. At $100\times100$, every classical
method succeeds on every scenario, while the learned policy succeeds on none
and is slower than A*. The defensible scope for this policy is therefore the
native small grid and closely related layouts, not scale-independent routing.

The multi-seed design changes how learned performance is interpreted. Scenario
confidence intervals from one frozen policy quantify test-case variation, not
training variance. Reporting seed-level values first ensures the uncertainty
shown for RL reflects independently learned policies. Likewise, the ablations
reveal interactions rather than a tidy component ranking. Double DQN and local
observations improve markedly over vanilla DQN and full-grid input, but
removing HER, shaping, or curriculum improves held-out-pair success from 96\%
to 100\% under this protocol. A one-factor intervention demonstrates
configuration sensitivity; in a nonlinear learning system it does not isolate
a context-free causal main effect.

\subsection{Limitations}
\label{sec:limitations}

The environment remains a two-dimensional discrete grid rather than a
continuous flight-dynamics simulator. Movement ignores vehicle turn radius,
altitude, acceleration, localization drift, communication constraints, and
multi-UAV interactions. Wind and energy are represented as spatial traversal
penalties rather than aerodynamic forces. The study therefore supports claims
about routing algorithms under controlled graph changes, not end-to-end flight
control.

The legacy contract permits departure from a cell that becomes blocked while
occupied. This avoids instantaneous failures caused by a toggle and is applied
equally to all methods, but it is only one possible physical interpretation.
Future work should compare escape, collision, and safety-buffer semantics.

Training uses one fixed static layout before dynamic fine-tuning. The benchmark
measures the resulting distribution shift explicitly, but does not establish
the performance of policies trained over randomized layout distributions.
The full-observation ablation is dimension-specific and cannot be applied
unchanged to larger grids, unlike the fixed-size local observation.

Only five independent policy seeds are available. Student-$t$ intervals expose
the resulting uncertainty, but remain wide for unstable variants and cannot
characterize rare optimization failures as precisely as a larger seed study.
The 310-scenario suite is broad within the grid model but is not a random sample
of all operational UAV environments.

Finally, timings describe these Python and PyTorch implementations on the
recorded Colab runtime class. Historical per-variant machine manifests were not
retained separately, so absolute cross-variant timing contrasts warrant extra
caution. The pipeline now prevents that provenance loss. Timings remain
auditable implementation-specific measurements, not hardware-independent
constants; operational conclusions require optimized code and embedded target
hardware.
